\section{Security Analysis}
\label{sec:security}

We analyze finite executions of the protocol in Section~\ref{sec:protocol}.
Each fixed four-member organization and the independent, equally weighted
four-validator committee have at most one statically Byzantine participant;
certificates require three distinct signatures. Honest participants share
legal genesis and a registry assigning each route one active configuration
and lock domain. Signing state is atomic and never forgotten or rolled
back; application transitions are deterministic and atomic, and followers
apply authenticated consecutive blocks with state and cursor updated
atomically. Claims concern states at these commit boundaries. We condition
on unforgeable signatures and binding canonical digests and original-part
commitments. Adapter output is treated as an untrusted candidate: a revision
takes effect only through a committed repair whose guards derive the
permitted slot, amount, and issuer and verify every non-target byte and the
complete block identity; installation accepts only bytes equal to the
committed Task body. Threshold adaptation supplies commitment-compatible
candidates, and its availability governs repair progress.
Delays are arbitrary. We first establish the original-execution interface,
then connect certification, reserve coverage, and repair. The consensus argument specializes the Tendermint locking interface~\cite{tendermint} to complete block identities.

\subsection{Original Execution and Certified Consumption}
\label{sec:security-identity}

A \emph{polka} is a same-round quorum of prevotes for a complete block identity or nil.

\begin{lemma}[Original-execution agreement]
\label{lem:security-original}
Under the locking and validity rules of the consensus interface, one original block body is admitted per
committed height, including its application context.
\end{lemma}
\begin{IEEEproof}
With original opening~1, a part leaf hashes its context-bound message
representative. Binding and the Merkle proof fix its bytes under the
original part-set header. Candidate bodies are reconstructed from their
parts or locally paired with parts; Locked and Valid transfers preserve
this pairing. Proposal, proof-of-lock (POL), lock, and commit comparisons use
$(H_b,P_b)$, and changing a download target clears the old body. A matching
current-round polka may trigger precommit before Proposal arrives; the
existing round and step guards still prevent a second precommit.

Same-round conflicting commits would require an honest validator to
precommit twice. For a commit of value $v$ in round $r$, take its two
honest signers and suppose a conflicting polka exists later in round
number. Choose the smallest such round $s>r$ and its earliest certificate
formation. Its quorum intersects those signers. That signer's vote
requires an earlier unlocking justification: a lower conflicting round
contradicts minimality, and an earlier same-round polka contradicts the
chosen formation event. Include nil polkas that authorize unlocking.
The argument uses signature events, not observation of the completed
commit certificate. Honest rounds advance monotonically, so a higher-round
vote by a commit signer follows its own earlier lock, even if the commit
certificate is assembled later. The validity of a vote cast for a locked value traces to earlier lock
rounds and ultimately to actual proposal validation. Finalize follows
a commit matching the candidate's complete block identity. Original
synchronization and replay preserve the body, height, and time, so
retries supply the same execution material.
\end{IEEEproof}

\begin{lemma}[Authenticated, unique consumption]
\label{lem:security-consumption}
A TXCer fixes its outputs and issuer. No two valid certificates
authorize conflicting consumption of the same $(x,i)$, and public
execution consumes that instance at most once.
\end{lemma}
\begin{IEEEproof}
Binding fixes the certified summary and owner-authorized transaction.
The fixed route sends competing requests to one organization. Two
three-member quorums intersect in at least two members, including an
honest member whose atomic input lock precedes signing. That member
cannot approve both conflicts. Budget release preserves the lock.
Final-input and certificate-input execution consult the same spend key;
fact replay produces neither another consumption nor another fee charge.
A recovered output has no final creation and its instance~0 stays
consumed, so no later claim on $(x,i)$ can succeed.
\end{IEEEproof}

\subsection{Covering Public and Hidden Commitments}
\label{sec:security-coverage}

For a complete CAL grant $g$, cumulative authorization $B_g$ gives each
member $s(B_g)=\lfloor2B_g/3\rfloor$ capacity. Let $C_g$ contain every
distinct fact whose three signatures have existed by the current prefix.
It includes unassembled and hidden certificates and never deletes
settled facts. The unique CAL allocation binds each fact to one
configuration, ResourceKey, and Grant ID. A fact $c$ has output cap
$a_c$, the sum of its CAL outputs. For $c\in C_g$, let $p_c$ be the
cumulative compensation paid for those outputs, $\mathsf{Rec}_c$ the part
repaid by source execution, and $z_c$ indicate that the payment of $c$
has executed successfully. The \emph{net compensation} is
$\bar p_c=p_c-\mathsf{Rec}_c$. Define
\begin{equation}
  w_c=\begin{cases}a_c,&z_c=0,\\\bar p_c,&z_c=1,\end{cases}
  \qquad Q_g=\sum_{c\in C_g}w_c.
  \label{eq:security-risk}
\end{equation}

\begin{lemma}[Source repayment]
\label{lem:security-recovery}
A compensated output leaves the Repaired state only inside the successful
execution of its source, in the same transition that sets $z_c=1$, and
the repayment is credited to the account that was debited. Hence
$0\le\mathsf{Rec}_c\le p_c\le a_c$, $\mathsf{Rec}_c=0$ while $z_c=0$,
and $\mathsf{Rec}_c=p_c$ once $z_c=1$. In particular $\bar p_c=p_c$ while
the source is pending and $\bar p_c=0$ afterwards, so no net compensation
survives the execution of the source.
\end{lemma}
\begin{IEEEproof}
Each registered promise leaves Open at most once, and compensation applies only to a funded Open
obligation, so $p_c\le a_c$. The only transition leaving Repaired is the
output-resolution step of a payment whose fact equals the obligation's
certificate fact and whose issuer equals the obligation's issuer; it
executes in the same overlay that records the payment as settled, so
$z_c$ and the repayment change together. That step resolves every output
of the source that has an obligation: Open becomes Fulfilled and Repaired
becomes Recovered. Outputs without an obligation incur no compensation.
Afterwards every promise of $c$ is closed, so no new compensation can
arise, and a repeated submission is a no-op. Compensation debits and
repayment credits both use the issuer account stored in the obligation,
which is the account of the CAL allocation of $c$.
\end{IEEEproof}

For any local approval, define $r_{m,c}=a_c-\mathsf{Applied}_{m,c}$,
where Applied is its cumulative CAL restoration.

\begin{lemma}[Residual capacity]
\label{lem:security-witness}
For every fact $c\in C_g$ and every honest member $m$ that approved $c$,
at every prefix,
$r_{m,c}\ge w_c$.
\end{lemma}
\begin{IEEEproof}
Until $m$ observes the successful execution of $c$'s payment, its CAL
restoration for $c$ is zero, so $r_{m,c}=a_c\ge w_c$ whatever the
compensation state. After it observes that execution, $m$ checks that the
reported repayment covers its recorded compensation and sets its recorded
repayment equal to it. The restoration target is $a_c$ minus a recorded
net compensation of zero, so $r_{m,c}=0$, while $w_c=\bar p_c=0$ by
Lemma~\ref{lem:security-recovery}. A member that has not yet observed its payment's public execution
retains $r_{m,c}=a_c\ge0=w_c$. The statement
covers a member that approved after following a compensation: it recorded
none and restores only its own debit.
\end{IEEEproof}

\begin{theorem}[Potential liabilities are funded]
\label{thm:security-coverage}
For every CAL grant, $Q_g\le B_g$. For each reserve account $x$, its total
open gap $D_x$ does not exceed its balance $A_x$.
\end{theorem}
\begin{IEEEproof}
Let $\mathcal A_{m,g}$ include all successful approvals under $g$, even
without a QC, and $L_{m,g}=\sum_{c\in\mathcal A_{m,g}}r_{m,c}$ across
Workers. Approval and restoration transfer capacity between Available
and Reserved. Top-ups add the difference in member shares, so
$L_{m,g}\le s(B_{m,g})\le s(B_g)$ for the locally observed authorization
$B_{m,g}$; grants only increase. Choose any fixed three honest members
$H_g$ of this organization. A QC has three distinct signers and at most
one member lies outside $H_g$, so the set $\mathsf{Sig}_c$ of members of
$H_g$ that signed $c$ has at least two elements. Honest members sign only
after approving, so Lemma~\ref{lem:security-witness} gives
$r_{m,c}\ge w_c$ for $m\in\mathsf{Sig}_c$, and
\begin{equation}
\begin{aligned}
  2Q_g
  &\le\sum_{m\in H_g}\sum_{\substack{c\in C_g\\m\in\mathsf{Sig}_c}}r_{m,c}\\
  &\le\sum_{m\in H_g}L_{m,g}
  \le3\lfloor2B_g/3\rfloor\le2B_g.
\end{aligned}
\label{eq:security-counting}
\end{equation}

Let $R_c$ and $d_c$ be registered remaining and discharged coverage,
with $p_c=d_c=R_c=0$ before registration. Registered records satisfy
$a_c=p_c+d_c+R_c$, and $z_c=1$ implies $R_c=0$. Therefore
$R_c\le w_c-\bar p_c$: while $z_c=0$ we have $\mathsf{Rec}_c=0$, so
$R_c\le a_c-p_c=w_c-\bar p_c$, and once $z_c=1$ both sides are zero.
Define $S_g=\sum\bar p_c$, $V_g=\sum R_c$, and $E_g=Q_g-S_g$, with sums
over $C_g$. Registration, compensation, repayment, and discharge
maintain public Usage.Spent$=S_g$ and Usage.Reserved$=V_g$.

Let $\mathcal G_x$ contain all grants backed by $x$. Their certified CAL
account equals the debited issuer account. Genesis establishes
$\sum_{g\in\mathcal G_x}(B_g-S_g)\le A_x$. External, non-protected
top-ups increase both sides equally. Compensation increases $S_g$ and
decreases $A_x$ by the same amount; by Lemma~\ref{lem:security-recovery},
repayment decreases $S_g$ and increases $A_x$ by the same amount. No
other protocol operation withdraws CAL reserves. Each Open obligation
corresponds to a distinct remaining promise, giving
\begin{equation}
\begin{aligned}
D_x&\le\sum_{g\in\mathcal G_x}V_g
    \le\sum_{g\in\mathcal G_x}E_g\\
   &\le\sum_{g\in\mathcal G_x}(B_g-S_g)\le A_x.
\end{aligned}
\label{eq:security-backing}
\end{equation}
Gross Paid and its recovery are each counted once; $E_g$ is a risk bound,
not available member capacity.
\end{IEEEproof}

\noindent\textbf{CAL admission consequence.}
At the execution prestate, including earlier successful block commands,
let $N_g$ be the distinct, valid, unregistered parent facts under $g$
required by matching, unconsumed missing inputs. They have $w_c=a_c$
and $p_c=R_c=0$. Every registered fact satisfies $R_c+\bar p_c\le w_c$,
and the facts are disjoint, so
\begin{equation}
 V_g+S_g+\sum_{c\in N_g}a_c\le Q_g\le B_g.
 \label{eq:security-admission}
\end{equation}
Repeated uses of one certificate contribute one term, not one term per
input. Sequential registration sees earlier reservations, and every
partial sum remains within this bound. First registration therefore
cannot fail solely for CAL authorization capacity, even when it precedes
a repayment in the same transition, since the repayment only lowers
$S_g$. Likewise, an Open amount $a\le D_x\le A_x$ has sufficient reserve
principal. Other resource and fee checks remain separate.

\subsection{Asset Conservation and Delay Neutrality}
\label{sec:security-conservation}

\begin{theorem}[CAL conservation]
\label{thm:security-conservation}
Let $U$ be the total value of final, unspent CAL instances, $A$ the total
CAL account balance, and $D$ the total value of Open gaps. Every
execution prefix satisfies
\begin{equation}
 U+A-D=C_0,
 \label{eq:security-conservation}
\end{equation}
where $C_0$ includes genesis outputs and all accounts, including top-up
sources.
\end{theorem}
\begin{IEEEproof}
A payment consumes final value $F$ and missing value $M$, producing
$F+M$. Let $J$ be its arriving outputs whose obligations were Open, and
$R$ those whose obligations were compensated. Outputs in $J$ were already
consumed, so they close gaps rather than add unspent value, and outputs
in $R$ receive no creation record. Hence $\Delta U=M-J-R$,
$\Delta D=M-J$, and the repayment gives $\Delta A=R$, so
$\Delta(U+A-D)=0$. Compensation of $a$ gives $\Delta A=\Delta D=-a$ and
$\Delta U=0$. Top-ups transfer between included accounts. Other
transitions preserve this projection.
\end{IEEEproof}

FUEL uses separate accounting. Each escrow maximum equals Held plus
cumulative Rewards, Burned, and Refunded. The global conserved sum
counts unspent FUEL, ordinary accounts, current Held, actual reward
balances, and cumulative burn. Unique stage and obligation identities
transfer value once; refunded outputs and credited rewards are not
counted again through cumulative escrow fields. This includes user-funded
and authorized organization-funded fees.

Conservation fixes totals but not who holds them. The next result shows
that delay does not change the final holders.

\begin{theorem}[Delay neutrality]
\label{thm:security-neutrality}
Let $\mathcal T$ be a finite set of payments such that (i) no two consume
the same instance or fee input, (ii) the relation ``consumes an output
created by'' is acyclic, and (iii) every certificate input of a payment in
$\mathcal T$ certifies an output of a payment in $\mathcal T$. Let
$\mathcal U$ be a finite set of authorized reserve increases. Start from a
clean genesis state with no coverage, promise, or obligation records and
zero CAL Reserved and Spent. Let $\rho_0$ execute $\mathcal T$ in an order
consistent with (ii), applying $\mathcal U$, and let $\rho$ be any
execution consisting of the payments of $\mathcal T$, the increases of
$\mathcal U$, and compensations of due Open obligations, in any order.
Suppose every payment and increase succeeds exactly once in both
executions. Then at the end of $\rho$ and of $\rho_0$:
\begin{enumerate}
\item the final, unspent CAL instances coincide, with the same owners and
amounts;
\item every reserve account holds the same CAL balance;
\item every CAL grant has the same authorization and
Reserved$=$Spent$=0$, every obligation is Fulfilled or Recovered, and
$\mathsf{Paid}_c=\mathsf{Rec}_c$ for every certificate.
\end{enumerate}
Fees and FUEL are not compared.
\end{theorem}
\begin{IEEEproof}
\emph{Holders.} By (i), both executions consume the same instances. In
$\rho$, an output consumed before its source executes receives an
obligation; when the source executes, the output is either Fulfilled,
with a creation record that is already consumed, or Recovered, with no
creation record. Either way it is consumed. Outputs consumed by no payment
of $\mathcal T$ are created as final and unspent in both executions.
Hence the unspent instances are the genesis and $\mathcal T$ outputs not
consumed by $\mathcal T$, in both.

\emph{Reserves.} For a reserve account $x$ let $\hat A_x$ be its balance
plus the amounts of Repaired obligations issued by $x$. Compensation
lowers $A_x$ and raises the Repaired total by the same amount; by
Lemma~\ref{lem:security-recovery}, repayment does the reverse on the
same account; creating or fulfilling an obligation changes neither term.
As in Theorem~\ref{thm:security-coverage}, only top-ups, compensation,
and repayment change reserve balances. Therefore
$\hat A_x=A_x(0)+\tau_x$ at every prefix, where $\tau_x$ is the total of
increases credited to $x$ so far. An obligation exists only for an output
consumed before its source executes. By (iii) that source executes in
$\rho$, and its execution leaves the obligation Fulfilled or Recovered.
So at the end no obligation is Open or Repaired, $A_x=A_x(0)+\tau_x$, and
this is also the value in $\rho_0$, where no compensation occurs.

\emph{Usage.} Every registered certificate is used by a missing input, so
by (iii) its source executes in $\rho$. That execution closes all of its
promises, so Remaining$=0$, and Lemma~\ref{lem:security-recovery} gives
$\mathsf{Paid}_c=\mathsf{Rec}_c$. Hence Reserved$=\sum$Remaining$=0$ and
Spent$=\sum(\mathsf{Paid}_c-\mathsf{Rec}_c)=0$. Authorizations change
only through $\mathcal U$.
\end{IEEEproof}

Intermediate states differ: while a compensated source is pending, its
issuer's balance is lower by the compensated amount, which
Theorem~\ref{thm:security-coverage} covers. The theorem is not a claim of
incentive compatibility, and it says nothing about a source that never
executes, whose compensated amount remains a net loss of its issuer.

\subsection{Authorized Repair and Representation Independence}
\label{sec:security-repair}

\begin{lemma}[Constrained repair authority]
\label{lem:security-repair}
Every newly effective repair item closes its unique, due,
execution-prestate Open obligation. A successful repair changes no
transaction bytes outside its authorized funding slots, even when
adaptation outputs are adversarial.
\end{lemma}
\begin{IEEEproof}
Let $N$ be the RSA modulus, $e$ its public exponent, and $r$ an opening in $\mathbb Z_N^*$. The domain-separated representative $\mathsf{FDH}(ctx,m)$ binds the message and its repair context. The relation
\begin{equation}
 C=\mathsf{FDH}(ctx,m)r^e=\mathsf{FDH}(ctx,m')r'^e\pmod N
\end{equation}
establishes
compatibility, not fresh approval; a published opening ratio can be
reused for the same message relation. Public guards independently derive
the slot, issuer, amount, and reserve-debit identity, check current base
and body digests, and verify all non-target bytes and complete block
commitments. Invalid candidates leave no changes. Exact successful
replays are no-ops; repackaged commands encounter the shared obligation
terminal state. Economic effects and the authorized revision Task
commit together. The adapter neither writes ledger state nor forges
committee authentication. The argument uses its verified output,
not evidence of fresh threshold participation.
\end{IEEEproof}

Let $\pi_{\mathrm{econ}}$ retain public assets, consumption, coverage,
obligations, grant usage, and fee states, omitting revision and
installation metadata.

\begin{lemma}[Batch equivalence and installation independence]
\label{lem:security-installation}
At its commit boundary, an accepted batch has the economic effect of
its ordered individual compensation steps starting from the same state,
at the same time, without intervening commands. For the same public
history, legal physical installation schedules have identical economic
projections.
\end{lemma}
\begin{IEEEproof}
Unique obligations and slots are canonically ordered. Each step reads
the shared overlay, including earlier account, Pending, and fee changes.
Failure discards all steps; success commits their aggregate effects with
one revision. Base is a revision number; Previous and Next authenticate
body digests. The new revision is Base$+1$.

Canonical reads committed logical revision $\rho_\ell$. An exact Task
authorizes body and parts at revision $r\le\rho_\ell$. Installation
advances if $\rho_p<r$, checks exact equality if $\rho_p=r$, and treats
$\rho_p>r$ as already covered without rollback. It changes no economic
variable. Queue progress crosses only completed entries: installing
$A_2$ at $A_1$ cannot skip intervening $B_1$. Therefore, for legal
schedules $\mu_1,\mu_2$ over one public prefix,
\begin{equation}
 \pi_{\mathrm{econ}}(S_{\mu_1})=\pi_{\mathrm{econ}}(S_{\mu_2}).
 \label{eq:security-installation}
\end{equation}
Batch and single-command histories can have different application hashes;
deterministic replay of the same original command history reproduces
its application hash at every height. Repairs preserve fixed
transaction fields, authorizations, output identities, successor references,
and complete block identities.

For a follower starting from the same local state, configuration, and
cursor, equal fixed payment fields and authenticated results induce
equal economic updates whenever both representations decode. Mutable
Funding/opening does not select Paid, Pending, or fee effects. Repair
commands are authenticated in full; execution-result commitments
authenticate their per-item effects and fee outputs. This separates
economic observation from independent verification of the latest mutable
historical body.
\end{IEEEproof}

\subsection{Protected Continuation and Progress}
\label{sec:security-composition}

Define $\mathsf{FreshRecv}(w,o,c)$ as authenticated receipt by an honest
owner who has not previously authorized consumption of that instance.
Freshness is a history predicate, not inferred from successful receipt
alone.

\begin{theorem}[Protected continuation]
\label{thm:security-continuation}
Following FreshRecv, the certified output cannot be reinterpreted or
consumed without owner authorization. A certified, publicly executed
successor consumes a final source or atomically establishes a funded
direct obligation. Source closure and repair preserve its established
principal effects without successor re-signing or re-execution.
\end{theorem}
\begin{IEEEproof}
Lemma~\ref{lem:security-consumption} fixes ownership and excludes conflicting
consumption. Atomic execution couples each missing consumption with its
obligation and outputs. Theorem~\ref{thm:security-coverage} supplies funding;
Theorem~\ref{thm:security-conservation} accounts for closure and late sources.
Lemma~\ref{lem:security-original} identifies the executed history, and
Lemmas~\ref{lem:security-repair}--\ref{lem:security-installation} restrict
later changes. Induction over the joint local/public transitions
composes these properties; INSTALL retains local consumption effects
but does not change public assets. Repair may change Pending, fees,
and reserve balances without repeating successor principal.
\end{IEEEproof}

Continuation requires responsive honest quorums, available evidence,
sufficient fees and Worker/work capacity, and fair delivery and inclusion.
Repair additionally requires advancing public time and available adaptation
contributions. Partial approvals and unused promises can retain capacity;
timeouts alone never authorize unlocking. These are progress conditions,
not a universal completion deadline.


\noindent\textbf{Data availability and guarantor risk.}
Three questions are separate. \emph{Safety:}
Theorem~\ref{thm:security-coverage} funds every certificate ever signed,
whether or not its source is submitted. \emph{Progress:} a source
can progress when an honest relay holds its complete submission material
(transaction, admission vector, organization authorization, direct input
certificates) and the remaining fee, capacity, quorum, and fair-scheduling conditions hold; delivery time is unbounded. A
deadline can therefore pass for ordinary reasons such as congestion or a
short outage. Compensation then commits when adaptation contributions
are available: the issuer advances the missing principal and is repaid
when the source executes.
When all sources of the payment set execute,
Theorem~\ref{thm:security-neutrality} shows final holders and reserves
unchanged. \emph{Economics:} the member approval endpoint authenticates
owner, budget, and inputs but not the caller's role, so any client can
act as collector. If that client never installs or submits the assembled
material and no other relay obtains it, the source stays unexecuted and,
once compensated, its issuer bears the amount as an unrecovered loss. Reserve sizing reflects the
expected rate of such sources; fees are not priced as insurance against
them.
